\textbf{Phần 1. Trắc nghiệm lựa chọn}
\setcounterpageref{ex}{0}

\begin{ex} %[4H2V09]
Lúc 7 giờ đúng, góc bé tạo bởi kim giờ và kim phút của đồng hồ là
\choice
{Góc nhọn có số đo $60^\circ$}
{\True Góc tù có số đo $150^\circ$}
{Góc tù có số đo $120^\circ$}
{Góc bẹt có số đo $180^\circ$}
\loigiai{
Lúc 7 giờ đúng, khoảng cách ngắn hơn giữa hai kim từ số 7 đến số 12 gồm 5 khoảng giờ.\\
Mỗi khoảng giờ tương ứng với $30^\circ$.\\
Số đo góc bé tạo bởi hai kim là:
$$5 \times 30^\circ = 150^\circ$$
Vì $90^\circ < 150^\circ < 180^\circ$ nên đây là góc tù có số đo $150^\circ$.\\
\textbf{Đáp án của câu hỏi là B}
}
\end{ex}

\begin{ex} %[4H2V09]
Cho hình vuông $ABCD$. Kẻ hai đường chéo $AC$ và $BD$ cắt nhau tại điểm $O$. Điểm $O$ là đỉnh của bao nhiêu góc vuông trong hình vẽ?
\choice
{1 góc vuông}
{2 góc vuông}
{3 góc vuông}
{\True 4 góc vuông}
\loigiai{
Hai đường chéo của hình vuông vuông góc với nhau tại giao điểm $O$.\\
Do đó tại đỉnh $O$ tạo thành 4 góc vuông gồm: $\widehat{AOB}, \widehat{BOC}, \widehat{COD}, \widehat{DOA}$.\\
\textbf{Đáp án của câu hỏi là D}
}
\end{ex}

\begin{ex} %[4H2V09]
Một góc tù có số đo gấp $3$ lần một góc nhọn. Biết góc tù đó lớn hơn góc nhọn là $80^\circ$. Số đo của góc tù đó là
\choice
{\True $120^\circ$}
{$100^\circ$}
{$135^\circ$}
{$150^\circ$}
\loigiai{
Coi số đo góc nhọn là 1 phần thì số đo góc tù là 3 phần bằng nhau như thế.\\
Hiệu số phần bằng nhau là:
$$3 - 1 = 2\text{ (phần)}$$
Số đo của góc nhọn là:
$$80^\circ : 2 = 40^\circ$$
Số đo của góc tù là:
$$40^\circ \times 3 = 120^\circ$$
\textbf{Đáp án của câu hỏi là A}
}
\end{ex}

\begin{ex} %[4H2V09]
Một hình tam giác có số đo góc thứ nhất là $30^\circ$, góc thứ hai có số đo gấp 3 lần góc thứ nhất. Tam giác đó là
\choice
{Tam giác có 3 góc nhọn}
{Tam giác có 1 góc tù}
{\True Tam giác vuông}
{Tam giác đều}
\loigiai{
Số đo của góc thứ hai là:
$$30^\circ \times 3 = 90^\circ$$
Vì tam giác có một góc bằng $90^\circ$ (góc vuông) nên tam giác đó là tam giác vuông.\\
\textbf{Đáp án của câu hỏi là C}
}
\end{ex}

\begin{ex} %[4H2V09]
Một người đứng quay một chiếc kim chỉ hướng bắt đầu từ hướng Đông, quay theo chiều kim đồng hồ một góc bằng một góc bẹt cộng thêm một góc vuông. Hỏi chiếc kim lúc này chỉ về hướng nào?
\choice
{Hướng Nam}
{\True Hướng Bắc}
{Hướng Tây}
{Hướng Đông}
\loigiai{
Từ hướng Đông quay theo chiều kim đồng hồ một góc bẹt ($180^\circ$) thì kim chỉ về hướng Tây.\\
Từ hướng Tây tiếp tục quay theo chiều kim đồng hồ một góc vuông ($90^\circ$) thì kim chỉ về hướng Bắc.\\
\textbf{Đáp án của câu hỏi là B}
}
\end{ex}

\begin{ex} %[4H2V09]
Từ 8 giờ đúng đến 8 giờ 45 phút cùng ngày, kim phút của đồng hồ đã quét được một góc có số đo là
\choice
{\True $270^\circ$}
{$180^\circ$}
{$135^\circ$}
{$90^\circ$}
\loigiai{
Từ phút 0 đến phút 45 là 45 phút, tương ứng với $\frac{45}{60} = \frac{3}{4}$ vòng quay của kim phút.\\
Góc kim phút đã quét được là:
$$360^\circ \times \frac{3}{4} = 270^\circ$$
\textbf{Đáp án của câu hỏi là A}
}
\end{ex}

\begin{ex} %[4H2V09]
Cho góc $\widehat{xOy} = 150^\circ$. Vẽ tia $Oz$ nằm giữa hai tia $Ox$ và $Oy$ sao cho góc $\widehat{xOz} = 90^\circ$. Góc $\widehat{zOy}$ là
\choice
{Góc tù có số đo $100^\circ$}
{Góc vuông có số đo $90^\circ$}
{Góc nhọn có số đo $45^\circ$}
{\True Góc nhọn có số đo $60^\circ$}
\loigiai{
Số đo của góc $\widehat{zOy}$ là:
$$\widehat{zOy} = 150^\circ - 90^\circ = 60^\circ$$
Vì $60^\circ < 90^\circ$ nên góc $\widehat{zOy}$ là góc nhọn có số đo $60^\circ$.\\
\textbf{Đáp án của câu hỏi là D}
}
\end{ex}

\begin{ex} %[4H2V09]
Lúc 11 giờ đúng, góc bé tạo bởi kim giờ và kim phút của đồng hồ có số đo là
\choice
{$60^\circ$}
{$45^\circ$}
{\True $30^\circ$}
{$15^\circ$}
\loigiai{
Khoảng cách ngắn hơn giữa hai kim lúc 11 giờ đúng gồm đúng 1 khoảng giờ (từ số 11 đến số 12).\\
Số đo góc tạo thành là:
$$1 \times 30^\circ = 30^\circ$$
\textbf{Đáp án của câu hỏi là C}
}
\end{ex}

\begin{ex} %[4H2V09]
Cho một góc bẹt đỉnh $O$. Người ta vẽ thêm $3$ tia chung gốc $O$ nằm cùng về một phía của đường thẳng chứa góc bẹt đó. Hỏi có tất cả bao nhiêu góc đỉnh $O$ được tạo thành?
\choice
{$8$ góc}
{\True $10$ góc}
{$12$ góc}
{$6$ góc}
\loigiai{
Góc bẹt gồm 2 tia đối nhau, cùng với 3 tia vẽ thêm thì tại đỉnh $O$ có tất cả:
$$2 + 3 = 5\text{ (tia chung gốc)}$$
Số góc tạo thành là:
$$5 \times (5 - 1) : 2 = 10\text{ (góc)}$$
\textbf{Đáp án của câu hỏi là B}
}
\end{ex}

\begin{ex} %[4H2V09]
Trong một hình tứ giác có hai góc vuông và một góc tù bằng $120^\circ$. Số đo của góc còn lại trong tứ giác đó là
\choice
{$90^\circ$}
{$75^\circ$}
{$45^\circ$}
{\True $60^\circ$}
\loigiai{
Tổng bốn góc trong hình tứ giác bằng $360^\circ$.\\
Số đo góc còn lại là:
$$360^\circ - (90^\circ + 90^\circ + 120^\circ) = 360^\circ - 300^\circ = 60^\circ$$
\textbf{Đáp án của câu hỏi là D}
}
\end{ex}

\begin{ex} %[4H2V09]
Lúc 7 giờ 30 phút, góc bé tạo bởi kim giờ và kim phút của đồng hồ là
\choice
{\True Góc nhọn có số đo $45^\circ$}
{Góc nhọn có số đo $30^\circ$}
{Góc vuông có số đo $90^\circ$}
{Góc tù có số đo $105^\circ$}
\loigiai{
Lúc 7 giờ 30 phút, kim phút chỉ vào số 6.\\
Kim giờ ở vị trí chính giữa số 7 và số 8 (tương ứng vị trí $7{,}5$).\\
Khoảng cách giữa hai kim là:
$$7{,}5 - 6 = 1{,}5\text{ (khoảng giờ)}$$
Số đo góc tạo thành là:
$$1{,}5 \times 30^\circ = 45^\circ$$
Vì $45^\circ < 90^\circ$ nên đây là góc nhọn có số đo $45^\circ$.\\
\textbf{Đáp án của câu hỏi là A}
}
\end{ex}

\begin{ex} %[4H2V09]
Cho góc nhọn $\widehat{AOB} = 40^\circ$. Vẽ góc $\widehat{BOC}$ kề với góc $\widehat{AOB}$ sao cho góc $\widehat{AOC}$ tạo thành một góc bẹt. Số đo của góc $\widehat{BOC}$ là
\choice
{$120^\circ$}
{$130^\circ$}
{\True $140^\circ$}
{$150^\circ$}
\loigiai{
Vì góc $\widehat{AOC}$ là góc bẹt nên $\widehat{AOC} = 180^\circ$.\\
Số đo của góc $\widehat{BOC}$ là:
$$\widehat{BOC} = 180^\circ - \widehat{AOB} = 180^\circ - 40^\circ = 140^\circ$$
\textbf{Đáp án của câu hỏi là C}
}
\end{ex}

\textbf{Phần 2. Trắc nghiệm đúng sai}
\setcounterpageref{ex}{0}

\begin{ex} %[4H2V09]
Xét tính đúng sai của các khẳng định sau về góc và hình học:
\choiceTF
{\True Góc nhọn lớn nhất có số đo tính theo độ nguyên là $89^\circ$.}
{\True Góc có số đo bằng hai lần góc bẹt tương ứng với một vòng quay trọn vẹn $360^\circ$.}
{Tổng số đo của hai góc nhọn luôn nhỏ hơn $90^\circ$.}
{\True Trong hình chữ nhật không phải hình vuông, hai đường chéo cắt nhau tạo thành 2 góc nhọn và 2 góc tù.}
\loigiai{
\begin{itemize}
    \item Số nguyên nhỏ hơn 90 lớn nhất là 89 nên góc nhọn lớn nhất có số đo độ nguyên là $89^\circ$. Khẳng định a Đúng.
    \item $2 \times 180^\circ = 360^\circ$ ứng với 1 vòng tròn trọn vẹn. Khẳng định b Đúng.
    \item Ví dụ $50^\circ + 50^\circ = 100^\circ > 90^\circ$ nên khẳng định hai góc nhọn cộng lại luôn nhỏ hơn $90^\circ$ là Sai.
    \item Hai đường chéo hình chữ nhật cắt nhau tạo thành 2 cặp góc đối đỉnh gồm 2 góc nhọn và 2 góc tù. Khẳng định d Đúng.
\end{itemize}
\begin{itemize}
    \item a) \textbf{Đúng}.
    \item b) \textbf{Đúng}.
    \item c) \textbf{Sai}.
    \item d) \textbf{Đúng}.
\end{itemize}
}
\end{ex}

\begin{ex} %[4H2V09]
Cho tam giác $ABC$ có góc $\widehat{A} = 90^\circ$, góc $\widehat{B} = 45^\circ$. Xét tính đúng sai của các khẳng định sau:
\choiceTF
{\True Tam giác $ABC$ là tam giác vuông tại đỉnh $A$.}
{\True Góc đỉnh $C$ có số đo bằng $45^\circ$.}
{\True Tam giác $ABC$ có 1 góc vuông và 2 góc nhọn bằng nhau.}
{Góc đỉnh $C$ là một góc tù.}
\loigiai{
\begin{itemize}
    \item Vì góc $\widehat{A} = 90^\circ$ nên tam giác $ABC$ vuông tại $A$. Khẳng định a Đúng.
    \item Số đo góc $\widehat{C}$ là: $180^\circ - (90^\circ + 45^\circ) = 45^\circ$. Khẳng định b Đúng.
    \item Tam giác có 1 góc vuông ($90^\circ$) và 2 góc nhọn bằng nhau ($45^\circ$). Khẳng định c Đúng.
    \item Vì $\widehat{C} = 45^\circ < 90^\circ$ nên là góc nhọn (không phải góc tù). Khẳng định d Sai.
\end{itemize}
\begin{itemize}
    \item a) \textbf{Đúng}.
    \item b) \textbf{Đúng}.
    \item c) \textbf{Đúng}.
    \item d) \textbf{Sai}.
\end{itemize}
}
\end{ex}

\textbf{Phần 3. Trả lời ngắn}
\setcounterpageref{ex}{0}

\begin{ex} %[4H2V09]
Cho góc bẹt $\widehat{xOy} = 180^\circ$. Kẻ tia $Ot$ nằm giữa hai tia $Ox$ và $Oy$ sao cho góc $\widehat{xOt} = 65^\circ$. Tính số đo của góc $\widehat{yOt}$ theo độ (chỉ ghi số).
\shortans{115}
\loigiai{
\textbf{Đáp số:} $115$\\
Số đo của góc $\widehat{yOt}$ là:
$$\widehat{yOt} = 180^\circ - 65^\circ = 115^\circ$$
}
\end{ex}

\begin{ex} %[4H2V09]
Một góc tù có số đo bằng $150^\circ$. Hỏi góc đó lớn hơn góc nhọn $30^\circ$ bao nhiêu độ?
\shortans{120}
\loigiai{
\textbf{Đáp số:} $120$\\
Số độ lớn hơn là:
$$150^\circ - 30^\circ = 120^\circ$$
}
\end{ex}

\begin{ex} %[4H2V09]
Lúc 1 giờ 30 phút, góc bé tạo bởi kim giờ và kim phút của đồng hồ có số đo bằng bao nhiêu độ?
\shortans{135}
\loigiai{
\textbf{Đáp số:} $135$\\
Lúc 1 giờ 30 phút, kim phút chỉ số 6, kim giờ ở vị trí chính giữa số 1 và 2 ($1{,}5$).\\
Khoảng cách giữa hai kim: $6 - 1{,}5 = 4{,}5$ khoảng giờ.\\
Số đo góc tạo thành là:
$$4{,}5 \times 30^\circ = 135^\circ$$
}
\end{ex}

\begin{ex} %[4H2V09]
Trong các hình: tam giác, tứ giác, ngũ giác, lục giác, hình nào có tổng số đo các góc trong bằng $360^\circ$? Hãy ghi số cạnh của hình đó.
\shortans{4}
\loigiai{
\textbf{Đáp số:} $4$\\
Hình tứ giác có tổng số đo 4 góc trong luôn bằng $360^\circ$. Số cạnh của hình tứ giác là 4.
}
\end{ex}

\begin{ex} %[4H2V09]
Cho $6$ tia chung gốc $O$. Hỏi có tất cả bao nhiêu góc đỉnh $O$ được tạo thành?
\shortans{15}
\loigiai{
\textbf{Đáp số:} $15$\\
Số góc tạo thành từ 6 tia chung gốc là:
$$6 \times (6 - 1) : 2 = 15\text{ (góc)}$$
}
\end{ex}

\begin{ex} %[4H2V09]
Một bạn học sinh vẽ liên tiếp 3 góc nhọn kề nhau có số đo lần lượt là $25^\circ, 35^\circ$ và $30^\circ$. Hỏi góc lớn nhất tạo bởi cạnh đầu tiên và cạnh cuối cùng có số đo bằng bao nhiêu độ?
\shortans{90}
\loigiai{
\textbf{Đáp số:} $90$\\
Số đo góc lớn nhất tạo thành là:
$$25^\circ + 35^\circ + 30^\circ = 90^\circ$$
}
\end{ex}
